% !PRINTED 504
\weberpage{504}

%<<CONT
metric group belongs to a group $Q$, then it belongs to $Q$ also within every other group $P$ of which $Q$ is a divisor.

The second part, however, follows from the fundamental properties of the \gesperrt{Galois} group.

To see this, we shall introduce the following notation. If any permutation $\pi$ of the group $P$ is applied to a function $\psi$ of the field $N$, let $\psi$ pass over into $\psi(\pi)$; if a permutation $\pi'$ is applied to $\psi(\pi)$, the result is the same as if the composed permutation $\pi\pi'$ had been applied to $\psi$, i.\,e.\ we have
\weq{6}{\psi(\pi)(\pi') = \psi(\pi\pi').}
If now $\psi$ remains unchanged by the permutation $\varkappa$, then
\[
\psi = \psi(\varkappa),
\]
and to this, by \S.\,149, a), any permutation $\pi$ from the group $P$ may be applied. From this it follows that
\[
\psi(\pi) = \psi(\varkappa\pi).
\]
If therefore $\psi(\pi)$ is also equal to $\psi$, then $\psi(\varkappa\pi) = \psi$ as well, i.\,e.\ the permutations which leave $\psi$ unchanged form a group, q.\,e.\,d.

If $\psi$ is a function in $N$ belonging to the group $Q$, so that
\weq{7}{\psi(\varkappa) = \psi(\varkappa_1) = \cdots = \psi(\varkappa_{q-1}) = \psi,}
and $\pi_1$ is any permutation of $P$ not belonging to $Q$, then we can apply $\pi_1$ to the equations (7) and obtain
\[
\psi(\varkappa\pi_1) = \psi(\varkappa_1\pi_1) = \cdots = \psi(\varkappa_{q-1}\pi_1) = \psi_1,
\]
and $\psi_1$ is different from $\psi$. If conversely, for some $\pi$,
\weq{8}{\psi_1 = \psi(\pi) = \psi(\varkappa\pi_1)}
holds, then $\pi$ must be contained in the sequence of substitutions
\weq{9}{Q\pi_1 = \varkappa\pi_1,\ \varkappa_1\pi_1,\ldots \varkappa_{\nu-1}\pi_1}
% [sic] The source uses the terminal index \nu-1, although Q is indexed through q-1 in (2) and (7); preserved source-faithfully.
for by applying $\pi^{-1}$ to (8) it follows that $\varkappa\pi_1\pi^{-1}$ must be contained in $Q$, say equal to $\varkappa_1$; but from this it follows that
\[
\pi = \varkappa_1^{-1}\varkappa\pi_1,
\]
which is indeed contained in (9).

We can accordingly state the following theorem:

4. \gesperrt{A function $\psi$ belonging to $Q$ passes, under all the permutations of a coset $Q\pi$, and}
%CONT>>
